for every @@M0@@ a locally free @@M1@@-Module of finite type @@M2@@ by putting:
@@M3@@
where @@M4@@ is the structural morphism @@M5@@.

We then see that it suffices to prove the proposition in the case where @@M6@@ is a projective bundle over @@M7@@, the case of a flag bundle following by induction. This case follows immediately from 1.1 (resp. 2.1).

\sgaStableTarget{sga6.e06.thm.3-2}\textbf{Theorem 3.2.} \uline{Let} @@M8@@ \uline{be a quasi-compact scheme}. \uline{The pre}-@@M9@@-\uline{ring} @@M10@@, \uline{defined from locally free} @@M11@@\uline{-Modules of finite type}, \uline{is a} @@M12@@-\uline{ring}~(V 2.4).

The homomorphism @@M13@@ being additive, it suffices, to check that it is a @@M14@@-homomorphism, to make the verification for elements of the form @@M15@@, where @@M16@@ is a locally free @@M17@@-Module of finite type. We are therefore reduced to showing that, for every pair of elements @@M18@@ of @@M19@@, classes of locally free @@M20@@-Modules of finite type @@M21@@, relations V 2.4.2 hold:
@@M22@@
@@M23@@

Since @@M24@@ is a quasi-compact scheme, we can find a finite covering of @@M25@@ by disjoint opens such that, on each, @@M26@@ and @@M27@@ have constant rank. The @@M28@@ of a finite direct sum of opens is the product of the @@M29@@ of these opens. It therefore suffices to check relations V 2.4.2 in each of the @@M30@@ of the opens of the covering. We can therefore assume @@M31@@ and @@M32@@ to have constant rank on @@M33@@.

